\begingroup
\setlength{\parskip}{4pt}
\setlength{\tabcolsep}{3pt}
\renewcommand{\arraystretch}{1.40}
\setstretch{1}
\footnotesize
\noindent\textbf{FICHA 01}\par
{\bfseries Ficha geométrica de evacuación longitudinal}\par
Corredor EMP.LI105 Pedregal--Censo--Cumbray, Simbal | Evaluación 2026\par
Sector asociado al pase de agua. Delimite hasta 50 m a cada lado y recorte en extremos de tramo o en el punto medio entre pases vecinos.\par
\begin{tabularx}{\linewidth}{|p{2.1cm}|p{1.1cm}|X|p{3.3cm}|}
\hline
\textbf{Código unidad} & \textbf{Tramo} & \textbf{Progresivas inicio y fin} & \textbf{Fecha y hora de medición}\\\hline
\rule{0pt}{4mm} & & & \\\hline
\end{tabularx}
Evaluador \rule{3.3cm}{0.25pt}\quad Pase de agua \rule{2.9cm}{0.25pt}\par
Coordenadas WGS84 \rule{2.4cm}{0.25pt}\quad Lluvia o mantenimiento \rule{2.2cm}{0.25pt}\par
Registre por margen los intervalos evaluados y los que tienen recorrido libre y continuo hasta una descarga. Cuente una sola vez los intervalos superpuestos. Anote barreras, acumulaciones y terminaciones cerradas con su evidencia.\par
\begin{tabularx}{\linewidth}{|p{1.15cm}|p{1.45cm}|p{1.4cm}|p{0.8cm}|p{0.8cm}|X|p{0.75cm}|}
\hline
\textbf{Margen} & \textbf{Prog. inicio} & \textbf{Prog. fin} & \textbf{Le (m)} & \textbf{Lc (m)} & \textbf{Descarga o interrupción} & \textbf{Foto}\\\hline
\rule{0pt}{4mm}&&&&&&\\\hline
\rule{0pt}{4mm}&&&&&&\\\hline
\rule{0pt}{4mm}&&&&&&\\\hline
\rule{0pt}{4mm}&&&&&&\\\hline
\rule{0pt}{4mm}&&&&&&\\\hline
\rule{0pt}{4mm}&&&&&&\\\hline
\end{tabularx}
\textbf{Consolidación del sector}\par
\begin{tabularx}{\linewidth}{|p{4.0cm}|X|X|X|}
\hline
\textbf{Margen} & \textbf{Evaluada Le (m)} & \textbf{Continua Lc (m)} & \textbf{Continuidad (\%)}\\\hline
Izquierdo &&&\\\hline
Derecho &&&\\\hline
Sector (suma de longitudes) &&&\\\hline
\end{tabularx}
CL (\%) = 100 $\times$ Lc total / Le total. Instrumentos auxiliares: wincha, GPS y cámara.\par
Observaciones \rule{9.0cm}{0.25pt}\par
Responsable \rule{3.8cm}{0.25pt}\qquad Firma \rule{3.1cm}{0.25pt}
\endgroup
